% Part: first-order-logic
% Chapter: sequent-calculus
% Section: rules-and-proofs

\documentclass[../../../include/open-logic-section]{subfiles}

\begin{document}

\iftag{FOL}
      {\olfileid[hi]{fol}{seq}{rul}}
      {\olfileid[hi]{pl}{seq}{rul}}

\olsection{नियम और \usetoken{p}{derivation}}

आगे $\Gamma, \Delta, \Pi, \Lambda$ को !!{sentence}s के परिमित क्रम मानें।

\begin{defn}[अनुक्रम]
\emph{अनुक्रम} इस रूप का व्यंजक है:
\[
\Gamma \Sequent \Delta
\]
जहाँ $\Gamma$ और $\Delta$, भाषा $\Lang L$ के !!{sentence}s के परिमित
(संभवतः रिक्त) क्रम हैं। $\Gamma$ को \emph{पूर्वपक्ष} और $\Delta$ को
\emph{उत्तरपक्ष} कहते हैं।
\end{defn}

\begin{explain}
अनुक्रम के पीछे सहज विचार यह है: यदि पूर्वपक्ष के सभी !!{sentence}s सत्य हों,
तो उत्तरपक्ष का कम-से-कम एक !!{sentence} सत्य हो। अर्थात् यदि
$\Gamma = \tuple{!A_1, \dots, !A_m}$ और
$\Delta = \tuple{!B_1, \dots, !B_n}$ हों, तो $\Gamma \Sequent \Delta$
तभी और केवल तभी सत्य है जब
\[
(!A_1 \land \cdots \land !A_m) \lif (!B_1 \lor \cdots \lor
!B_n)
\]
सत्य हो। दो विशेष दशाएँ हैं: $\Gamma$ रिक्त हो या $\Delta$~रिक्त हो। जब
$\Gamma$ रिक्त हो, अर्थात् $m = 0$, तब $\quad
\Sequent \Delta$ तभी और केवल तभी सत्य है जब $!B_1 \lor \dots \lor !B_n$
सत्य हो। जब $\Delta$ रिक्त हो, अर्थात् $n = 0$, तब $\Gamma \Sequent \quad$
तभी और केवल तभी सत्य है जब $\lnot(!A_1 \land \dots \land !A_m)$ सत्य हो।
अनुक्रम को वैध तभी और केवल तभी कहते हैं जब उससे संबद्ध !!{sentence} वैध हो।
\end{explain}

यदि $\Gamma$, !!{sentence}s का क्रम है, तो $\Gamma, !A$ उस परिणाम को दर्शाता
है जो $!A$ को~$\Gamma$ के दाएँ सिरे पर जोड़ने से मिलता है (और $!A,
\Gamma$ उस परिणाम को, जो $!A$ को~$\Gamma$ के बाएँ सिरे पर जोड़ने से मिलता
है)। यदि $\Delta$ भी !!{sentence}s का क्रम है, तो $\Gamma,
\Delta$ दोनों क्रमों का संयोजन है।

\begin{defn}[प्रारंभिक अनुक्रम]
\emph{प्रारंभिक अनुक्रम} ऐसा अनुक्रम है जो
\iftag{prvFalse,prvTrue}{इनमें से किसी एक रूप का हो:
  \begin{enumerate}
    \item $!A \Sequent !A$
    \tagitem{prvTrue}{$\quad \Sequent \ltrue$}{}
    \tagitem{prvFalse}{$\lfalse \Sequent \quad$}{}
  \end{enumerate}}
{रूप $!A \Sequent !A$ का हो} जहाँ $!A$ भाषा का कोई भी !!{sentence} है।
\end{defn}

अनुक्रम-कलन की !!^{derivation}s अनुक्रमों के विशेष वृक्ष हैं। उनके शीर्षस्थ
अनुक्रम प्रारंभिक अनुक्रम होते हैं; और यदि कोई अनुक्रम एक या दो अन्य अनुक्रमों
के नीचे हो, तो उसे अनुमान-नियम से ठीक प्रकार निकलना चाहिए। $\Log{LK}$ के नियम
दो मुख्य प्रकारों में बँटते हैं: \emph{तार्किक} और \emph{संरचनात्मक} नियम।
तार्किक नियम का नाम उस !!{main
  operator} पर होता है जो निचले अनुक्रम में $!A$ और/या $!B$ वाले !!{sentence}
का है। हर नियम के दो रूप हैं: एक में वह !!{sentence}, जिसमें !!{operator} है,
बाईं ओर और दूसरे में वही !!{sentence} दाईं ओर निष्कर्षित होता है।

\end{document}
